\documentclass[11pt]{article}
\usepackage[a4paper,margin=27mm]{geometry}
\usepackage[T1]{fontenc}
\usepackage{lmodern,microtype,amsmath,amssymb,amsthm,mathtools,booktabs,array,longtable}
\usepackage{xcolor}
\usepackage[unicode,colorlinks=true,linkcolor=black,citecolor=black,urlcolor=blue!45!black]{hyperref}
\usepackage{enumitem}
\setlist{nosep}
\hypersetup{pdftitle={Prime-Power Obstructions to Flagged CRT Transport and an Exact Divisor-Gap Reduction for abc},pdfauthor={ChatGPT},pdfsubject={Research supplement; no unconditional abc proof claimed}}
\setlength{\parskip}{4pt}
\setlength{\parindent}{1.2em}
\newtheorem{theorem}{Theorem}[section]
\newtheorem{proposition}[theorem]{Proposition}
\newtheorem{lemma}[theorem]{Lemma}
\newtheorem{corollary}[theorem]{Corollary}
\theoremstyle{definition}\newtheorem{definition}[theorem]{Definition}
\theoremstyle{remark}\newtheorem{remark}[theorem]{Remark}
\newcommand{\rad}{\operatorname{rad}}
\newcommand{\dist}{\operatorname{dist}}
\newcommand{\FCRT}{\mathrm{FCRT}}
\newcommand{\SCRT}{\mathrm{SCRT}}
\newcommand{\PBT}{\mathrm{PBT}}
\newcommand{\pos}[1]{\left(#1\right)_+}
\newcommand{\N}{\mathbb N}
\newcommand{\Z}{\mathbb Z}
\newcommand{\R}{\mathbb R}
\title{Prime-Power Obstructions to Flagged CRT Transport\\and an Exact Divisor-Gap Reduction for \(abc\)}
\author{ChatGPT}
\date{5 September 2026}
\begin{document}
\maketitle
\begin{center}\small
Research manuscript and reproducible supplement.\\
Ordinary proofs and exact computations; the Lean companion has not been compiled.\\
No unconditional proof or disproof of the \(abc\) conjecture is claimed.
\end{center}
\begin{abstract}
We continue the FCRT-1 route of the ABCConjecture repository at commit
\texttt{6118955}. We prove an arithmetic obstruction that is independent of
pigeonhole estimates: if \(p^e\Vert c\), \(M=p^e\ge4\), and a nonempty proper
prime-power packet of \(c-1\) is congruent to \(-1\pmod M\), then
\(c\ge M(2M-1)\). The threshold is attained for infinitely many moduli.
Consequently, balanced endpoints \(c=2^e3^f\) have no proper flags at either
source. An elementary approximation argument produces infinitely many such
endpoints with \(49\mid c-1\), hence \(c>\rad(c(c-1))\). This refutes a
universal flag-existence shortcut even on positive-defect endpoints, whose
absolute defects can moreover be arbitrarily large; it does
not refute FCRT-1's uniform boundary conjecture or \(abc\).
On this stratum we identify the entire FCRT optimum with an exclusive
partition optimum and express it exactly as scalar radical defect plus a
one-dimensional divisor-gap penalty. We also give a multiplicative arithmetic
certificate for general FCRT configurations, an exact minimum-product deletion
normal form, and a finite dynamic program with a proved invariant. A
25-digit endpoint is fully factored and attains the scalar optimum despite
having no flags. All global uniform estimates remain explicit open obligations.
\end{abstract}
\bigskip

\section{Scope, provenance, and endpoint notation}
The source snapshot is
\begin{center}\small\texttt{6118955d20b4edd32e577e06d1060f3945358dd9}.\end{center}
The September 4 endpoint, anchored-prefix, and independent-audit notes
\cite{repoendpoint,repoanchor,repoaudit} are the immediate antecedents.
We retain their mathematical FCRT-1 definition, rather than replacing its
admissibility condition with a desired height inequality. In particular, a
flag must have an actual nonempty proper prime-power face. Results below are
proved here; priority or novelty beyond comparison with those notes is not
asserted. Computational evidence is separated from universal statements.

Let \(a,b,c\) be positive integers, \(a+b=c\), and \(\gcd(a,b)=1\).
Put \(R=\rad(abc)\). For each \(p\mid c\), let \(e_p=v_p(c)\), and set
\[
 I=\{p:e_p\ge2\},\quad d_p=p^{e_p-1},\quad x_p=\log d_p,
 \qquad J=\{q:q\mid ab\}.
\]
For \(S\subseteq I\), \(T\subseteq J\), write
\[
 D_S=\prod_{p\in S}d_p,\quad m_S=\prod_{p\in S}p^{e_p},\quad
 Q_T=\prod_{q\in T}q,
\]
\[
 a_T=\prod_{q\in T,\ q\mid a}q^{v_q(a)},\qquad
 b_T=\prod_{q\in T,\ q\mid b}q^{v_q(b)}.
\]
Empty products are one. These are \emph{whole-prime-power packets}: an
exponent may not be split between faces. A block \((S,T)\) has nonempty
source and sink sets and is admissible and saturated when
\begin{equation}\label{eq:block}
 m_S\mid a_T+b_T,\qquad D_S\le Q_T.
\end{equation}
The selected blocks have pairwise disjoint source sets and pairwise disjoint
sink sets. Let \(I_0,J_0\) be the respective residual sets. A residual sink
has zero or one owner in \(I_0\). A selected block may emit zero or one token
to one \(p\in I_0\), witnessed by
\begin{equation}\label{eq:flag}
 \varnothing\ne U\subsetneq T,\qquad p^{e_p}\mid a_U+b_U.
\end{equation}
The logarithmic credit is
\[
 \rho=\min\{\log Q_T-\log D_S,\log Q_U\}.
\]
No credit is re-emitted. Several independent tokens may have the same target.
These rules are precisely the one-hop, capped FCRT-1 rules of the snapshot.

Let \(X=\sum_{p\in I}x_p\), \(Y=\log Q_J\), \(h=\log c\),
\(r=\log R\). Unique factorization and pairwise coprimality give
\begin{equation}\label{eq:height}
 \prod_{p\in I}d_p=\frac{c}{\rad(c)},\qquad
 X-Y=h-r,\qquad \Delta=\max\{h-r,0\}.
\end{equation}
We reserve \(B_{\FCRT},B_{\SCRT},B_{\PBT}\) for \emph{logarithmic}
optimized boundaries. Their exponentials are called boundary factors.

\section{A concrete multiplicative certificate}
This section supplies a mathematical arithmetic-to-accounting construction.
It is not a claim that the repository's Lean owner types have now been
inhabited or connected by compiled code.

For a used block \(\nu\), define the rational number
\begin{equation}\label{eq:creditfactor}
 f_\nu=\min\{Q_{T_\nu}/D_{S_\nu},Q_{U_\nu}\};
\end{equation}
for an unused block put \(f_\nu=1\). Saturation implies \(f_\nu\ge1\).
For a residual source \(p\), let
\[
 K_p=\prod_{q\in J_0:\ o(q)=p}q
       \prod_{\nu:\ t(\nu)=p}f_\nu,
 \qquad
 F=\prod_{p\in I_0}\max\{d_p/K_p,1\}.
\]
Everything in this definition uses positive integers, rational arithmetic,
finite sets, and divisibility. Its logarithm is exactly the original clipped
FCRT boundary.

\begin{theorem}[Arithmetic once-charge certificate]\label{thm:certificate}
Every configuration satisfying the actual disjointness, saturation, owner,
and face conditions above satisfies
\begin{equation}\label{eq:rationalbridge} c\le R F.\end{equation}
Consequently \(\Delta\le\log F\). The construction does not assume this
inequality as a field of its input.
\end{theorem}
\begin{proof}
For every residual source,
\(d_p\le K_p\max\{d_p/K_p,1\}\).
Multiplication, disjointness, and at-most-once ownership imply
\begin{align*}
 \prod_{p\in I}d_p
 &\le \Bigl(\prod_\nu D_{S_\nu}\Bigr)
       \Bigl(\prod_{p\in I_0}K_p\Bigr)F\\
 &\le \Bigl(\prod_\nu D_{S_\nu}f_\nu\Bigr) Q_{J_0} F
 \le \Bigl(\prod_\nu Q_{T_\nu}\Bigr)Q_{J_0}F=Q_JF.
\end{align*}
Unused sinks only increase the middle upper bound because their primes are
at least two. Substitute \eqref{eq:height} and multiply by \(\rad(c)\).
Since \(F\ge1\), taking logarithms also bounds the positive part.
\end{proof}

\begin{remark}[What the certificate does and does not use]
Congruence and witness-face conditions make the configuration arithmetically
admissible. The inequality itself uses disjointness and the surplus cap.
This distinction is essential: dropping arithmetic conditions preserves a
valid inequality but may destroy all additional arithmetic information.
The construction maps source weights to \(\log d_p\), block weights to
\(\log D_S,\log Q_T\), witness weights to the \emph{actual} \(\log Q_U\),
and owners to the given partial maps. The finite partitions explicitly
establish the source and sink mass decompositions required by the earlier
independent audit \cite{repoaudit}.
\end{remark}

\section{Optimal faces as minimum-product deletions}
Fix a source \(p\), and put \(M=p^{e_p}\ge4\).
All primes in \(J\) are units modulo \(M\). Define
\[
 g_q=\begin{cases}q^{v_q(a)}&q\mid a,\\q^{-v_q(b)}&q\mid b\end{cases}
 \quad\text{in }(\Z/M\Z)^\times,
 \qquad \Phi(T)=\prod_{q\in T}g_q.
\]
Then \(\Phi(U)=-1\) is equivalent to \(M\mid a_U+b_U\), and
\(\Phi(J)=-1\).
For any nonempty \(T\subseteq J\), define
\begin{equation}\label{eq:z}
 z_p(T)=\min\{Q_D:\varnothing\ne D\subseteq T,
                         \ \Phi(D)=-\Phi(T)\},
\end{equation}
with value \(+\infty\) when the set is empty.

\begin{theorem}[Exact optimal-face normal form]\label{thm:normal}
For a fixed saturated block \((S,T)\) and fixed residual target \(p\), a
proper flag exists if and only if \(z_p(T)<\infty\). If one exists, the
maximum credit factor over its legal witnesses is
\begin{equation}\label{eq:optimalfactor}
 f^*_{p,S,T}=\frac{Q_T}{\max\{D_S,z_p(T)\}}.
\end{equation}
In particular its discarded logarithmic surplus is
\begin{equation}\label{eq:loss}
 \log\frac{Q_T}{D_S}-\log f^*_{p,S,T}
       =\max\{\log z_p(T)-\log D_S,0\}.
\end{equation}
\end{theorem}
\begin{proof}
For complementary subsets \(U,D\) in \(T\), multiplication in the unit
group gives
\(\Phi(U)=-1\) if and only if \(\Phi(D)=-\Phi(T)\).
The nonempty condition on \(D\) says \(U\ne T\). Also \(D=T\) cannot
satisfy its displayed label: it would imply \(\Phi(T)=-\Phi(T)\), hence
\(M\mid2\), contrary to \(M\ge4\). Thus \(U\) is automatically nonempty.
Since \(Q_U=Q_T/Q_D\), maximizing the face weight is equivalent to minimizing
\(Q_D\). Monotonicity of the minimum in \eqref{eq:creditfactor} gives
\eqref{eq:optimalfactor}; cancellation gives \eqref{eq:loss}.
\end{proof}

When \(\Phi(T)=-1\), the deletion label in \eqref{eq:z} is one. For a
general block it is \emph{not} necessarily one. This distinction prevents
an incorrect application of an anchored identity-deletion result to an
arbitrary cross-source block. Optimizing each fixed target independently is
valid: increasing its incoming credit never increases the clipped boundary.

\subsection{Exact finite computation}
Process the distinct primes \(q_1,\ldots,q_n\) of \(T\). A table indexed
by reachable residues stores the least product of a \emph{nonempty} subset
and a witness. Initially it is empty. At stage \(k\), retain all previous
entries, insert the singleton \((g_{q_k},q_k)\), and, using only the previous
table, insert
\[
      (r g_{q_k},wq_k)\quad\text{from every previous entry }(r,w).
\]
At each residue retain the smaller product. The empty subset is not allowed
to overwrite a nonempty identity entry.

\begin{proposition}[Dynamic-program invariant]
After stage \(k\), each present table entry is the minimum product among
nonempty subsets of \(\{q_1,\ldots,q_k\}\) with that residue; an absent
entry means no such subset exists. Thus the entry at \(-\Phi(T)\) computes
\eqref{eq:z} exactly.
\end{proposition}
\begin{proof}
A nonempty subset either omits \(q_k\), is the singleton \(\{q_k\}\), or
is \(\{q_k\}\) together with a nonempty earlier subset. These are exactly
the three transitions. Multiplication by the positive integer \(q_k\)
preserves order, so discarding a larger product with the same earlier residue
cannot remove a later optimum. Induction proves the assertion.
\end{proof}
The sparse table has at most the size of the subgroup generated by the
labels. Arithmetic-operation complexity is \(O(n|H|)\) in a dense
implementation; this excludes bit-complexity and makes no claim of polynomial
time in the bit length of \(M\). No numerical logarithm is used.

\subsection{An arithmetic cost obstruction}
\begin{proposition}\label{prop:cost}
Suppose \(\Phi(T)=-1\), and let \(D\) be a nonempty identity deletion.
Write \(E=\max_{q\in T}v_q(ab)\) and
\(V_T=a_Tb_T/Q_T\). Then
\[
 Q_D^E\ge M+1,\qquad Q_DV_T\ge M+1.
\]
Consequently, full-surplus reuse at a block requires both
\(D_S^E\ge M+1\) and \(D_SV_T\ge M+1\).
\end{proposition}
\begin{proof}
The positive coprime integers \(a_D,b_D\) are congruent modulo \(M\).
They are unequal: equality and coprimality would make both one, contrary to
a nonempty packet. Hence \(|a_D-b_D|\ge M\) and
\(a_Db_D\ge\max\{a_D,b_D\}\ge M+1\).
Furthermore \(a_Db_D\le Q_D^E\), and
\(a_Db_D=Q_DV_D\le Q_DV_T\).
Full reuse means \(Q_D\le D_S\) for some deletion, by
Theorem~\ref{thm:normal}. Substitute this inequality.
\end{proof}

\section{A sharp obstruction for unit-gap endpoints}
The next result is stronger than a weight obstruction: in its range the
required face does not exist at any weight.

\begin{lemma}[Bilinear factor gap]\label{lem:gap}
Let \(M\ge2\), \(s,t\ge1\) be integers, and suppose
\[
 C+t=s(Mt+1),\qquad C\ne M.
\]
Then \(C\ge2M-1\).
\end{lemma}
\begin{proof}
If \(s=t=1\), the equation gives \(C=M\), excluded by hypothesis.
If \(t\ge2\), then
\[
 C=s(Mt+1)-t\ge (M-1)t+1\ge2M-1.
\]
If \(t=1\), necessarily \(s\ge2\), and
\(C=s(M+1)-1\ge2M+1\).
\end{proof}

\begin{theorem}[Prime-power no-face threshold]\label{thm:noface}
Let \(c\) be a positive integer and \(p^e\Vert c\), with \(e\ge2\).
Put \(M=p^e\). If a nonempty proper prime-power packet \(U\) of \(c-1\)
satisfies \(M\mid1+b_U\), then
\begin{equation}\label{eq:threshold} c\ge M(2M-1).\end{equation}
In fact the conclusion holds for any factorization
\(c-1=du\), \(d,u>1\), with \(u\equiv-1\pmod M\).
\end{theorem}
\begin{proof}
Write \(c=MC\); exactness of \(p^e\Vert c\) implies \(p\nmid C\),
so \(C\ne M\). The complementary packet gives
\(c-1=du\) with \(d,u>1\). Since \(du\equiv-1\pmod M\) and
\(u\equiv-1\pmod M\), we have \(d\equiv1\pmod M\). Thus
\[
 d=Mt+1,\qquad u=Ms-1\qquad(s,t\ge1).
\]
Expanding \((Mt+1)(Ms-1)=MC-1\) and dividing by \(M\) gives
\(C+t=s(Mt+1)\). Lemma~\ref{lem:gap} yields \eqref{eq:threshold}.
\end{proof}

The exact prime-power assumption is significant. Without it, the case
\(C=M\) is possible; the weaker threshold \(c\ge M^2\) is all the same
argument yields. The proof does not rely on estimating the number of
prime factors or on random residue heuristics.

\begin{proposition}[Sharpness on actual prime-power packets]
For \(M=2^{2k+1}\), \(k\ge1\), put \(c=M(2M-1)\). Then
\(v_2(c)=2k+1\) and a proper packet modulo the full \(2\)-power exists.
\end{proposition}
\begin{proof}
We have \(c-1=(M-1)(2M+1)\). Since \(M\equiv2\pmod3\),
\(\gcd(M-1,2M+1)=\gcd(M-1,3)=1\). Therefore the two factors are
unions of disjoint whole-prime-power packets. Both exceed one, and
\(M-1\equiv-1\), \(2M+1\equiv1\pmod M\).
The factor \(2M-1\) is odd, proving the valuation assertion.
\end{proof}

\begin{corollary}[Balanced two-source endpoints]\label{cor:balanced}
Suppose \(c=2^e3^f\), \(e,f\ge2\), and
\begin{equation}\label{eq:balanced}
 \max\{2^e,3^f\}<2\min\{2^e,3^f\}-1.
\end{equation}
For \((1,c-1,c)\), no nonempty proper packet of the full sink set is
compatible with either source. In particular, \emph{no} FCRT-1 configuration
at this endpoint contains a used flag, regardless of its selected blocks.
\end{corollary}
\begin{proof}
For either source modulus \(M\), the complementary factor \(C=c/M\)
is less than \(2M-1\), by \eqref{eq:balanced}. Apply
Theorem~\ref{thm:noface}. A proper face of any block is also a proper face
of the full sink set, so choosing a smaller block cannot evade the result.
\end{proof}

\section{Infinitely many positive-defect endpoints without flags}
\begin{theorem}\label{thm:infinite}
There are infinitely many distinct primitive triples \((1,c-1,c)\) such
that \(c=2^e3^f\), \eqref{eq:balanced} holds, \(49\mid c-1\), and
\begin{equation}\label{eq:hit}
 \rad(c(c-1))\le\frac{6(c-1)}7<c.
\end{equation}
All their proper FCRT flags are absent.
\end{theorem}
\begin{proof}
The number \(\alpha=2\log3/\log2\) is irrational. Otherwise an equality
between a positive power of 2 and a positive power of 3 would contradict
unique factorization. The elementary pigeonhole approximation argument
produces positive integers \(u_k,v_k\), with \(v_k\to\infty\), such that
\(u_k-\alpha v_k\to0\). For completeness, place the \(Q+1\) fractional
parts of \(0,\alpha,\ldots,Q\alpha\) into \(Q\) equal intervals. A
collision gives \(1\le v\le Q\) and an integer \(u\) with
\(|v\alpha-u|\le1/Q\). These denominators cannot remain bounded as
\(Q\to\infty\), since irrationality gives a positive minimum error on any
fixed finite set. Pass to a subsequence with unbounded denominators; its
numerators are eventually positive.

Set \(e_k=21u_k\), \(f_k=42v_k\). Then
\[
 \log\frac{2^{e_k}}{3^{f_k}}
       =21\log2\,(u_k-\alpha v_k)\longrightarrow0.
\]
Both factors tend to infinity. Their larger-to-smaller ratio tends to one,
so \eqref{eq:balanced} holds eventually. The elementary congruences
\(2^{21}\equiv1\pmod{49}\) and \(3^{42}\equiv1\pmod{49}\)
imply \(49\mid2^{e_k}3^{f_k}-1\).

A number divisible by \(7^2\) has radical at most one seventh of itself.
Also \(\rad(c)=6\) and \(\gcd(c,c-1)=1\). This proves
\eqref{eq:hit}. Select an increasing subsequence of \(c\) to ensure
distinctness, and invoke Corollary~\ref{cor:balanced}.
\end{proof}


\begin{corollary}[Unbounded absolute defect without flags]\label{cor:unbounded}
For every integer \(k\ge2\), there are infinitely many balanced two-source
endpoints without flags such that \(7^k\mid c-1\) and
\[
 \frac{c}{\rad(c(c-1))}>\frac{7^{k-1}}6.
\]
Consequently the scalar defect \(\Delta\) is unbounded on the no-flag
stratum. For every fixed threshold there are infinitely many no-flag
endpoints with defect above that threshold.
\end{corollary}
\begin{proof}
The congruences
\[
 2^{3\cdot7^{k-1}}\equiv1\pmod{7^k},\qquad
 3^{6\cdot7^{k-1}}\equiv1\pmod{7^k}
\]
follow by induction from \(2^3\equiv3^6\equiv1\pmod7\): raising a number
congruent to one modulo \(7^j\) to the seventh power makes it congruent to
one modulo \(7^{j+1}\), as the binomial expansion shows. In the proof of
Theorem~\ref{thm:infinite}, replace the constants \(21,42\) by
\(3\cdot7^{k-1},6\cdot7^{k-1}\), keeping \(k\) fixed during the
approximation argument. The same ratio tends to one, and
\(\rad(c-1)\le(c-1)/7^{k-1}\). Now let \(k\) increase to exceed any
prescribed absolute defect threshold. There is still no estimate of this
defect as a positive proportion of \(\log R\).
\end{proof}

\begin{remark}[Exact statement refuted]
Theorem~\ref{thm:infinite} refutes the following proposed shortcut, including
its ``all sufficiently large'' variant: every positive-defect two-source
unit-gap endpoint has at least one proper compatible face. Corollary~\ref{cor:unbounded}
also rules out rescue by imposing only a sufficiently large absolute defect. It also prevents
a global application of any sufficient condition that would force such a
face at \emph{every} endpoint in this stratum. It does not contradict the
anchored-prefix theorem \cite{repoanchor}; that theorem has additional
hypotheses, which cannot all hold here. No original FCRT uniform gate is
refuted. Nor does a lower bound \(c/R>7/6\) refute \(abc\): it does not
establish that \(c/R^{1+\varepsilon}\) is unbounded for any fixed
\(\varepsilon>0\).
\end{remark}

\section{The remaining finite problem is exactly a divisor gap}
Let the hypotheses of Corollary~\ref{cor:balanced} hold. Write
\[
 Q=\rad(c-1),\quad A=2^{e-1},\quad B=3^{f-1},\quad D=\log(AB/Q).
\]
In this section assume \(AB>Q\), so \(D=\Delta>0\).

\begin{theorem}[Exact reduction on the no-flag stratum]\label{thm:reduction}
At these endpoints,
\begin{equation}\label{eq:allopt}
 B_{\FCRT}=B_{\SCRT}=B_{\PBT}
 =\log\min_{H\mid Q}\left(\max\{A/H,1\}\max\{BH/Q,1\}\right).
\end{equation}
Equivalently, for the nonempty real interval
\(\mathcal I=[\log(Q/B),\log A]\),
\begin{equation}\label{eq:distance}
 B_{\FCRT}=D+\min_{H\mid Q}\dist(\log H,\mathcal I).
\end{equation}
In particular, the scalar lower bound is attained if and only if there is
a divisor \(H\mid Q\) with \(Q/B\le H\le A\).
\end{theorem}
\begin{proof}
Every compatible nonempty block has the full sink set: a proper one would
contradict Corollary~\ref{cor:balanced}. Disjointness permits at most one
such block. The block with both sources is not saturated, because \(AB>Q\).
A saturated singleton-source block with all sinks has exactly the same
boundary as assigning every sink exclusively to that source. Thus it cannot
improve upon the exclusive partition optimum. Flags do not exist, so all
three optima coincide.

An unused residual sink can be assigned without increasing the boundary.
Since \(Q\) is squarefree, assigning all sinks between the two sources is
the same as choosing a divisor \(H\mid Q\); the capacities are \(H,Q/H\).
This proves \eqref{eq:allopt}. To obtain \eqref{eq:distance}, put
\(t=\log H\), \(L=\log(Q/B)\), \(U=\log A\). Since \(L<U\), direct
consideration of \(t<L\), \(L\le t\le U\), and \(t>U\) gives
\[
 (\log A-t)_++(\log B-\log Q+t)_+
       =D+\dist(t,[L,U]).
\]
Minimize over the finite divisor set.
\end{proof}

\subsection{A dense-divisor upper bound, not an \(abc\) estimate}
The maximum ratio of consecutive divisors is a classical object in the
study of dense divisors \cite{weingartner}. We include the elementary
squarefree calculation to make the precise inequality self-contained.
Write \(Q=q_1\cdots q_n\), with distinct increasing primes, and define
\[
 G(Q)=\max_{1\le j\le n}\frac{q_j}{q_1\cdots q_{j-1}}.
\]
The empty denominator is one, so \(G(Q)\ge2\).

\begin{lemma}\label{lem:dense}
The largest ratio of consecutive positive divisors of squarefree \(Q\)
is exactly \(G(Q)\).
\end{lemma}
\begin{proof}
Adjoin the primes in increasing order. In logarithmic coordinates, the new
divisor set is the union of the old set and its translate by \(\log q_j\).
If their convex hulls are disjoint, the new bridge gap is
\(\log(q_j/(q_1\cdots q_{j-1}))\); all other gaps are old ones. If the
hulls overlap, a gap in the union cannot cross either hull endpoint in its
interior, since those endpoints are themselves divisor logarithms. It therefore
lies in one hull and between two consecutive points of that translated or
untranslated set. Thus no gap exceeds the previous maximum.
This gives the upper bound by induction.
For the reverse bound, whenever \(q_j>q_1\cdots q_{j-1}\), a divisor
smaller than \(q_j\) uses only earlier primes and is at most their full
product. Hence that product and \(q_j\) are consecutive divisors of the
full \(Q\). The case \(j=1\) supplies the initial gap from one.
\end{proof}

\begin{corollary}\label{cor:densebound}
On the positive-defect no-flag stratum,
\[
 D\le B_{\FCRT}\le D+\frac12\max\{\log G(Q)-D,0\}.
\]
If \(\log G(Q)\le D\), the scalar optimum is attained.
\end{corollary}
\begin{proof}
The interval \(\mathcal I\) has length \(D\) and intersects
\([0,\log Q]\), since \(A,B\ge1\). If it contains a divisor logarithm,
the penalty is zero. Otherwise it lies strictly inside one consecutive-divisor
gap, of length \(g\le\log G(Q)\). The nearer endpoint is at distance at
most \((g-D)/2\) from the interval. Apply Theorem~\ref{thm:reduction}.
\end{proof}

There is no proved uniform bound here for either \(D\) or the endpoint's
divisor-gap penalty. A theorem about the density of integers with dense
divisors does not automatically apply pointwise to the special sequence
\(\rad(2^e3^f-1)\). The exact reduction changes the finite optimization
problem; it does not eliminate the Diophantine difficulty.

\section{An exact 25-digit endpoint and computational replay}
Take
\[
 e=41,\quad f=26,\quad
 M=2199023255552,\quad N=2541865828329,
\]
\[
 c=5589622068988418728132608.
\]
Direct integer arithmetic gives
\begin{align*}
 c-1&=7^2\cdot439\cdot857\cdot2729\cdot292183\cdot380261663,\\
 Q&=7\cdot439\cdot857\cdot2729\cdot292183\cdot380261663.
\end{align*}
All six displayed factors are prime; the companion program checks this by
deterministic trial division. The two no-face inequalities are
\[
 N<2M-1,\qquad M<2N-1.
\]
Consequently neither source has a proper flag, while \(R=6Q=6(c-1)/7<c\).
Nevertheless the divisor
\[
 H=2729\cdot380261663=1037734078327
\]
attains the scalar optimum. Indeed
\[
 H\le A=1099511627776,\qquad
 Q/H=769481753663\le B=847288609443.
\]
These are exact integer comparisons. Thus
\[
 \exp B_{\FCRT}=\frac{AB}{Q}
  =\frac{931603678164736454688768}{798517438426916961161801}
  =\frac{c}{R}.
\]
The absence of flags therefore does \emph{not} imply failure of the full
uniform-boundary strategy. It identifies a stratum requiring a different
mechanism, here an ordinary sink partition.

\subsection{Reproducibility protocol}
The standard-library Python companion performs: independent dynamic-program
versus exhaustive-subset comparisons; bounded direct checks of the no-face
threshold; exact rational FCRT optimization for the repository's finite
witnesses; deterministic primality and factorization checks for the endpoint
above; and arithmetic checks on a large symbolic member with exponents
\((399,252)\). In this replay there were 12,231 ordered primitive triples with \(c\le200\),
26,832 dynamic-program/exhaustive-table comparisons, 5,082 arithmetic-cost
checks, and 3,700 concrete owner-bridge checks. The separate unit-gap scan
through \(c=200000\) checked 6,741 eligible endpoint--modulus pairs and
35,756 proper unitary packets, with no counterexample. The dense-divisor
identity was checked on 12,159 squarefree values through 20,000; the bilinear
factor lemma was checked on 297,381 nonexceptional bounded parameter triples.
Its JSON output records exact bounds, factors, and selected configurations. All assertions are executable checks rather than
floating-point estimates. The results and invocation are in the accompanying
verification report. Bounded tests do not prove a uniform theorem, and a
Python replay is not a Lean kernel proof.

\section{Independent route and source audit}
\subsection{IUT: retain the route without assuming its comparison}
The original IUT III discusses the three indeterminacies, log-Kummer
compatibility, and log-volume comparison in the proof of Corollary 3.12
\cite{iut}. Inspection of those passages does not supply the missing
source-faithful global construction in the repository. Its own audit still
distinguishes local accounting interfaces from actual arithmetic data.
Joshi's report \cite{joshi} takes a distinct position: it criticizes both
sides of the controversy and claims that his arithmetic Teichmuller spaces
provide additional structures. These are attributed claims, not results
validated by this manuscript. We neither import the disputed comparison as
an axiom nor declare the entire IUT route false on the basis of a simplified
model. No IUT bridge is closed here.

\subsection{Anchored selection and recent literature}
The anchored code-fibre lemma in Bae's revised manuscript
\cite{bae} is an elementary sufficient selection principle. Its proof
selects two indices in the same fibre with a permitted positive difference.
It does not establish the arithmetic reservoir hypotheses at every \(abc\)
endpoint. Theorem~\ref{thm:infinite} exhibits infinitely many endpoints
where any sufficient condition implying a proper flag must fail.

Lichtman's discussion of the almost-all \(abc\) result \cite{lichtman}
concerns quantitative exceptional-set bounds; such a bound does not itself
exclude infinitely many exceptions. Dujella's prime-power tuple theorem
\cite{dujella} controls a different Diophantine configuration. No uniform
map from the endpoint packets above into that theorem's hypotheses has been
constructed here. These sources are relevant background, not an imported
solution to the uniform gate.

\subsection{Other maintained routes}
The repository's Pell, Mersenne/Wieferich, synchronized-packet, Steinberg,
and Frey/Szpiro ledgers identify separate unresolved valuation or uniformity
obligations. This manuscript does not settle those obligations and does not
retire those routes. In particular, fixed-parameter finiteness, polynomial
function-field analogies, or identities for radical defects cannot be promoted
to a constant uniform in all integer triples without another proof. The
present ordinary results use neither GRH nor \(abc\), IUT, a conjectural
prime-producing polynomial, or a conjecture about Wieferich primes.

\section{Formalization boundary and remaining uniform target}
The mathematical proofs above precede the Lean companion. That companion
contains candidate proof scripts for the bilinear arithmetic obstruction,
the proper-divisor implication, the real interval identity, and concrete
integer certificates. It has \emph{not} been compiled in this execution
environment. Accordingly there is no claimed Lean declaration count, no
claimed successful axiom audit, and no claimed unconditional closed term
of type \texttt{ABCConjecture}. The infinite-family approximation argument,
the full arithmetic owner construction, the dynamic-program invariant, and
the complete optimization theorems are not fully formalized by that file.

The pinned repository toolchain is \texttt{leanprover/lean4:v4.32.0}.
Official release notes for 4.32.1 and 4.32.2 describe kernel soundness fixes
\cite{lean321,lean322}; the latter also warns that the affected checking path
included comparator. This is not evidence that the repository exploited a
bug or that any particular theorem is false. It does mean that the eventual
trusted replay should rebuild compatible dependencies on a patched toolchain,
run a fresh-process check, and, where appropriate, an updated independent
checker. Merely printing a familiar list of axioms is not a substitute for
that check.

For the original target, a sufficient FCRT statement remains
\begin{equation}\label{eq:uniform}
 \forall\varepsilon>0\ \exists C_\varepsilon\in\R\
 \forall (a,b,c),\qquad
 B_{\FCRT}(a,b,c)\le\varepsilon\log\rad(abc)+C_\varepsilon.
\end{equation}
Together with Theorem~\ref{thm:certificate}, it gives the usual
\(c\le K_\varepsilon\rad(abc)^{1+\varepsilon}\).
Neither \eqref{eq:uniform} nor its negation is established.
On the positive-defect no-flag stratum, Theorem~\ref{thm:reduction} says
precisely that the boundary is the sum of two nonnegative terms:
\[
 \log\frac{c}{\rad(c(c-1))}
 \quad\text{and}\quad
 \min_{H\mid\rad(c-1)}\dist\!\left(\log H,
       [\log(Q/B),\log A]\right).
\]
A uniform sublinear bound for their sum is equivalent to such bounds for
both terms separately: necessity follows from nonnegativity, and sufficiency
uses \(\varepsilon/2\) in each bound. Thus further work must control
arithmetic radical defect as well as any transport fragmentation. Creating a
formal hypothesis with those bounds would not constitute proving them.

\paragraph{Repository integration status.}
This supplement is based on the pinned commit and is provided as an additive
patch. No remote branch, commit, pull request, or merge was created in this
session. Existing build imports and verified-status ledgers are deliberately
unchanged. A local research result, a locally replayed computation, a compiled
Lean proof, and a merged remote theorem are four different deliverables.

\begin{thebibliography}{20}\small
\bibitem{repoendpoint} ChatGPT, \emph{Endpoint residue cubes and proper-subface flagged CRT surplus}, ABCConjecture repository, 4 Sept. 2026, pinned commit \texttt{6118955}; research file \texttt{ABC\_ENDPOINT\_RESIDUE\_CUBE\_FLAGGED\_CRT\_2026\_09\_04.md}.
\bibitem{repoanchor} ChatGPT, \emph{Anchored-prefix flagged CRT}, same repository and commit; research file \texttt{ABC\_ANCHORED\_PREFIX\_FLAGGED\_CRT\_2026\_09\_04.md}.
\bibitem{repoaudit} \emph{Independent Lean audit of flagged CRT surplus}, same repository and commit; research file \texttt{ABC\_FLAGGED\_CRT\_LEAN\_INDEPENDENT\_AUDIT\_2026\_09\_04.md}.
\bibitem{weingartner} A. Weingartner, \emph{Practical numbers and the distribution of divisors}, Quarterly J. Math. 66 (2015), 743--758. \href{https://arxiv.org/abs/1405.2585}{arXiv:1405.2585}.
\bibitem{iut} S. Mochizuki, \emph{Inter-universal Teichmuller theory III: Canonical splittings of the log-theta-lattice}, author-hosted manuscript, especially Corollary 3.12 and Remark 3.11.1. \href{https://www.kurims.kyoto-u.ac.jp/~motizuki/Inter-universal\%20Teichmuller\%20Theory\%20III.pdf}{Author's PDF}, inspected 5 Sept. 2026.
\bibitem{joshi} K. Joshi, \emph{Final Report on the Mochizuki-Scholze-Stix Controversy}, \href{https://arxiv.org/abs/2505.10568}{arXiv:2505.10568v1} (2025).
\bibitem{bae} J. H. Bae, \emph{Unbounded logarithmic limsup in Erdos Problem 684 via shifted carry scheduling}, \href{https://arxiv.org/abs/2604.23784}{arXiv:2604.23784v3}, revised 3 Sept. 2026, Lemma 5.1. Earlier versions are not used.
\bibitem{lichtman} J. D. Lichtman, \emph{The abc conjecture is true almost always}, \href{https://arxiv.org/abs/2505.13991}{arXiv:2505.13991} (2025).
\bibitem{dujella} A. Dujella, \emph{Prime-power Diophantine tuples}, \href{https://arxiv.org/abs/2609.03448}{arXiv:2609.03448v1}, 3 Sept. 2026.
\bibitem{lean321} Lean development team, \emph{Lean 4.32.1 release notes}, 22 July 2026. \url{https://lean-lang.org/doc/reference/latest/releases/v4.32.1/}.
\bibitem{lean322} Lean development team, \emph{Lean 4.32.2 release notes}, 28 July 2026. \url{https://lean-lang.org/doc/reference/latest/releases/v4.32.2/}.
\end{thebibliography}
\end{document}
